\chapter{Conclusion and Discussion}
\label{chap:conclusion-discussion}

\section{Summary of Achievements}
\label{section:summary-of-achievements}

\begin{figure}[h]
    \centering
    \fbox{\parbox{0.9\textwidth}{\centering\vspace{1em}
    \textit{[To be completed at project conclusion. Will summarize delivered features, evaluation results, and degree to which the problem statement was addressed.]}
    \vspace{1em}}}
    \caption{Summary placeholder}
    \label{fig:summary-placeholder}
\end{figure}

\section{Known Limitations}
\label{section:limitations}

\begin{itemize}[leftmargin=80pt]
    \item \textbf{TCAS data currency:} TCAS admission requirements change each academic year. The system requires manual data updates at the start of each TCAS cycle and does not automatically detect changes.

    \item \textbf{CLO-level mapping:} The system uses program-level PLO data (มคอ.2) only. Course-level CLO data (มคอ.3) is not ingested, as features requiring CLO mapping (e.g., skill progress tracking for enrolled students) are out of scope for this project.

    \item \textbf{O*NET Thai career mapping:} O*NET occupational data is in English and uses US career taxonomy. Manual mapping to Thai career titles and KU graduate employment patterns is required and introduces approximation.

    \item \textbf{Portfolio scorer:} Designed and documented but not implemented in the MVP due to data dependency on structured per-faculty portfolio rubrics from the admissions office.

    \item \textbf{Evaluation ground truth:} The MRR evaluation test set is manually curated with limited size (30--50 profiles), which constrains the statistical reliability of the recommendation quality metrics.

    \item \textbf{Interest profile validity:} The interest discovery interface uses a five-step adaptive elicitation process as a lightweight proxy for established vocational interest instruments such as Holland's Self-Directed Search. While the five-step design is substantially more reliable than a two-question approach, it remains shorter than a full validated instrument (42+ items for the Self-Directed Search). Meta-analytic research on validated interest inventories reports approximately 50\% accuracy in predicting eventual career choice even with full assessments; the accuracy of KUru's abbreviated elicitation is expected to be lower. The system is designed accordingly --- recommendations are framed as a starting point for exploration rather than a definitive match, with the $\alpha$ blending mechanism progressively shifting recommendation weight toward behavioural signals as the student interacts with the system.

    \item \textbf{Behavioural blending cold start:} The $\alpha$ blending mechanism that combines RIASEC fit scores with behavioural fit scores provides no benefit until sufficient interaction data has accumulated. For early users of the system, collaborative signal is sparse or absent, and the system falls back toward pure RIASEC recommendations. The fallback hierarchy is: (1) sufficient similar-profile users exist $\rightarrow$ collaborative signal applied normally; (2) few similar users $\rightarrow$ $\alpha$ is held higher regardless of individual interaction count; (3) no similar users $\rightarrow$ pure RIASEC ($\alpha = 1.0$). Logged-in users always receive better personalisation than guests, as guest interaction signals are session-scoped and do not persist.
\end{itemize}

\section{Future Work}
\label{section:future-work}

\begin{itemize}[leftmargin=80pt]
    \item \textbf{KU enrolled student features (Phase 2):} Once the pre-admission system is validated, extend scope to serve currently enrolled KU students. This would require มคอ.3 CLO ingestion, a course enrollment interface, and a skill gap dashboard. The current knowledge graph schema is designed to support this extension without architectural changes.

    \item \textbf{Portfolio scorer:} Build the AI-assisted portfolio evaluation feature once structured admission rubrics are available from the KU admissions office.

    \item \textbf{Automated TCAS data refresh:} Monitor mytcas.com for changes each TCAS cycle and automatically notify the data administrator when updates are detected, triggering a review and re-ingestion cycle with KU faculty.

    \item \textbf{Thai career taxonomy alignment:} Develop a systematic mapping between O*NET occupations and THACO (Thai Standard Occupational Classification) to improve career recommendation accuracy for Thai students.

    \item \textbf{True collaborative filtering (Phase 2):} The current MVP implements single-user behavioural blending --- the $\alpha$ mechanism in UC-07 weights each student's \textit{own} interaction history against their initial RIASEC scores. True collaborative filtering --- using interaction patterns of students with similar RIASEC profiles to inform each other's recommendations --- requires a critical mass of user data and is deferred to Phase~2. The $\alpha$ integration point in Stage~6 of the recommendation pipeline is designed to accommodate this extension without architectural changes. A prerequisite is robust handling of \textit{profile pollution}: a poor initial RIASEC vector clusters a student incorrectly, degrading the collaborative signal for all students in that cluster. The five-step adaptive elicitation in UC-01 is designed to reduce this risk before true collaborative filtering is introduced.

    \item \textbf{Longitudinal evaluation:} Track recommendation quality by following up with admitted students after one year to assess whether their actual program experience matched the system's prediction.
\end{itemize}
